\documentclass[12pt,reqno]{amsart}
\usepackage{/Users/johnmyers/GoogleDrive/tex/handout,/Users/johnmyers/GoogleDrive/tex/customcommands,polynom}
\usepackage{caption}

\hdr{MAT 240}{\S16.5 {\small Integrals in cylindrical and spherical coordinates, part 1}}

\begin{document}

\textbf{Suggested problems:} 1-41 odd

\bigskip
\hrule

\bigskip
\prob Using the diagram below, write down equations that relate the cylindrical coordinates $(r,\theta,z)$ of a point to its cartesian (rectangular) coordinates $(x,y,z)$.

\medskip

\includegraphics[scale=1.25]{2.pdf}

\medskip
\textcolor{red}{From cylindrical to cartesian:
	\[x = r\cos{(\theta)}, \quad y = r\sin{(\theta)}, \quad z=z.
	\]
From cartesian to cylindrical:
	\[r^2 = x^2+y^2, \quad \tan{(\theta)} = \frac{y}{x}, \quad z=z.
	\]}

\bigskip
\prob

\begin{enumerate}
\item Plot the point with cylindrical coordinates $(2,2\pi/3,1)$ and find its cartesian coordinates.

\bigskip
\textcolor{red}{See class notes for plot. Cartesian coordinates $= (-1,\sqrt{3},1)$.}
\bigskip

\item Find the cylindrical coordinates of the point with cartesian coordinates $(3,-3,-7)$.

\bigskip
\textcolor{red}{We have $(\sqrt{18},-\pi/4,-7)$.}
\bigskip

\end{enumerate}

\prob

\begin{enumerate}
\item Describe the surface whose equation in cylindrical coordinates is $z=r$. \textcolor{red}{See class notes.}
\item Describe the surface whose equation in cylindrical coordinates is $r=1$. \textcolor{red}{See class notes.}
\item Describe the surface whose equation in cylindrical coordinates is $\theta =\pi/4$. \textcolor{red}{See class notes.}
\end{enumerate}

\prob Describe in cylindrical coordinates a wedge of cheese cut from a cylinder $4$ cm high and $6$ cm in radius; this wedge subtends an angle of $\pi/6$ at the center.

\bigskip
\textcolor{red}{It is described by the inequalities $0 \leq z \leq 4$, $0 \leq r \leq 6$, and $0 \leq \theta \leq \pi/6$.}
\bigskip

\newpage
\prob Using the diagram below, express the volume element $dV$ in cylindrical coordinates.

\medskip
\includegraphics[scale=1.25]{pic1.pdf}

\bigskip
\textcolor{red}{We have $dV = r drdzd\theta$.}
\bigskip

\prob For each object below, set up a triple integral of a function $f$ over the object.

\begin{enumerate}
\item \textcolor{white}{blank} \\
\includegraphics[scale=1.25]{pic3.pdf} \textcolor{red}{$\int_0^1 \int_0^3 \int_0^5 f(x,y,z) \ dzdydz$}
\item \textcolor{white}{blank} \\
\includegraphics[scale=1.25]{pic4.pdf} \textcolor{red}{$\int_0^{2\pi} \int_0^4 \int_0^1 f(r,\theta,z) \ r dzdrd\theta$}
\item \textcolor{white}{blank} \\
\includegraphics[scale=1.25]{pic5.pdf} \textcolor{red}{$\int_0^{\pi/2} \int_0^2 \int_0^4 f(r,\theta,z) \ r dzdrd\theta$}
\item \textcolor{white}{blank} \\
\includegraphics[scale=1.25]{pic7.pdf} \textcolor{red}{$\int_0^{\pi} \int_0^4 \int_0^r f(r,\theta,z) \ r dzdrd\theta$}
\end{enumerate}

\bigskip
\prob Evaluate
	\[\int_W \sin{(x^2+y^2)} \ dV
	\]
where $W$ is the solid cylinder with height 4 and with base of radius 1 centered on the $z$-axis at $z=-1$.

\bigskip
\textcolor{red}{
\begin{align*}
\int_0^{2\pi} \int_0^1 \int_{-1}^3 \sin{(r^2)}r \ dz dr d\theta  &= 4 \int_0^{2\pi} \int_0^1 \sin{(r^2)} r \ dr d\theta \\
&= -2 \int_0^{2\pi} ( \cos{(1)} -1 )  \ d\theta \\
&= -4\pi(\cos{(1)} -1 ).
\end{align*}}
\bigskip

\prob Find the volume between the cone $z=\sqrt{x^2+y^2}$ and the plane $z=10+x$ above the disk $x^2+y^2 \leq 1$.

\bigskip
\textcolor{red}{
\begin{align*}
\int_0^{2\pi} \int_0^1 \int_r^{10+r\cos{(\theta)}} r \ dzdrd\theta &= \int_0^{2\pi} \int_0^1 ( 10r+r^2\cos{(\theta)} -r^2 ) \ drd\theta \\
&= \int_0^{2\pi} \left( \frac{14}{3} + \frac{1}{3} \cos{(\theta)} \right) \ d\theta \\
&= \frac{28\pi}{3}.
\end{align*}}



\end{document}